\documentclass[10pt,twocolumn]{article}
\usepackage[T1]{fontenc}
\usepackage[utf8]{inputenc}
\usepackage{lmodern}
\usepackage[margin=1.8cm,top=2.4cm,bottom=2cm,columnsep=0.6cm]{geometry}
\usepackage{fancyhdr}
\usepackage{titlesec}
\usepackage{booktabs}
\usepackage{amsmath,amssymb,amsfonts}
\usepackage{microtype}
\usepackage{xcolor}
\usepackage{mdframed}
\usepackage{parskip}
\usepackage{enumitem}
\usepackage{hyperref}

\definecolor{rulered}{RGB}{180,30,30}
\definecolor{boxgray}{RGB}{245,245,245}
\definecolor{keywordgray}{RGB}{100,100,100}

\pagestyle{fancy}
\fancyhf{}
\fancyhead[L]{\textsc{\small Claude Mag}}
\fancyhead[C]{\small\textit{Machine Learning Research Digest}}
\fancyhead[R]{\small 2026-10-05}
\fancyfoot[C]{\small\thepage}
\renewcommand{\headrulewidth}{0.8pt}

\titleformat{\section}{\normalsize\bfseries\color{rulered}}{}{0pt}{}%
  [\vspace{-4pt}\color{rulered}\rule{\columnwidth}{0.4pt}]
\titlespacing{\section}{0pt}{8pt}{4pt}

\hypersetup{colorlinks=true,urlcolor=rulered,linkcolor=black}

\begin{document}
\setlength{\parindent}{0pt}
\setlength{\parskip}{4pt}

{\small\color{keywordgray}\textbf{Keywords:}
  diffusion models \textbullet{}
  inference-time scaling \textbullet{}
  Feynman-Kac correction \textbullet{}
  molecular free energy}\par\vspace{4pt}

{\LARGE\bfseries\raggedright
  Error-Corrected Inference-Time Scaling for Imperfect Diffusion Models}\par\vspace{4pt}

{\large\itshape\raggedright
  The Energy-based Feynman-Kac Corrector Removes Systematic Model Bias
  at Inference Time, Achieving Accurate Molecular Free-Energy Profiles
  Where Standard Particle Methods Fail}\par\vspace{6pt}

{\small\textbf{By} Zuokai Wen, Louis Grenioux, Weinan E \& Jiequn Han
  (Princeton University; Flatiron Institute, Simons Foundation)}\par\vspace{2pt}

{\small\color{keywordgray}arXiv:2610.01933 \textbullet{} October 2026}
\par\vspace{2pt}\noindent{\color{rulered}\rule{\columnwidth}{0.6pt}}\vspace{6pt}

Inference-time scaling for diffusion models---running more particles through sequential
Monte Carlo---reduces variance but cannot remove the systematic bias introduced by an
imperfect pretrained score function. The Energy-based Feynman-Kac Corrector (EBFKC)
derives dynamics that track a prescribed probability path exactly in the population limit
regardless of model error, and demonstrates accurate molecular free-energy profiling
on alanine dipeptide and tetrapeptide where standard methods fail.

\begin{mdframed}[backgroundcolor=boxgray,linecolor=rulered,linewidth=1pt,
    innerleftmargin=6pt,innerrightmargin=6pt,
    innertopmargin=6pt,innerbottommargin=6pt]
\textbf{\small AT A GLANCE}\par\vspace{4pt}
\begin{description}[leftmargin=0pt,labelsep=4pt,itemsep=2pt,parsep=0pt,
    font=\small\bfseries]
  \item[Problem:] Standard inference-time scaling inherits the systematic bias of an
    imperfect diffusion model; more particles reduce variance but not bias.
  \item[Why it matters:] Scientific applications (molecule generation, free-energy
    computation) require accurate tail distributions, not just good mode coverage.
  \item[Method:] Feynman-Kac potential derived from an energy-based reference corrects
    the imperfect score on the fly; approximated via sequential Monte Carlo with
    variance-controlling guidance.
  \item[Result:] Closely matches ground-truth free-energy surfaces on alanine dipeptide
    and tetrapeptide; standard SMC baselines retain several $k_BT$ of error.
  \item[Evidence:] Gaussian mixture models, particle systems, two alanine peptides under
    annealing and reward tilting.
\end{description}
\end{mdframed}

\section{The Big Picture}

Inference-time scaling has emerged as a key technique for improving diffusion model
quality: by running a sequential Monte Carlo (SMC) procedure with many particles, one
can steer generation toward a target distribution or reward signal without retraining
the model. Applications range from image generation with aesthetic rewards to
molecular conformation sampling with energy-based targets.

All existing approaches, however, share a silent assumption: that the pretrained
diffusion model's score function is exact. In practice, score functions are trained
on finite data with finite compute and carry systematic approximation error.
Standard SMC reduces the statistical error of the particle approximation but cannot
remove this model-level bias: every particle evolves under the same imperfect dynamics,
so the distribution they jointly represent converges to the wrong target.

\section{Why Existing Approaches Fall Short}

\textbf{Monte Carlo only reduces variance, not bias.} Adding more particles in standard
SMC reduces $O(1/\sqrt{N})$ variance, leaving a fixed systematic error proportional
to the score function mismatch $\epsilon_t = \hat{s}_\theta - \nabla\log p_t$.

\textbf{Path tracking error.} Imperfect scores cause particles to drift off the
prescribed probability path, accumulating errors that compound over diffusion timesteps.

\textbf{Reward tilting.} Reward-guided diffusion reweights samples assuming a correct
model; systematic score error introduces additional bias into the tilted distribution
that no reweighting scheme can correct post-hoc.

\section{Core Mechanism}

\textbf{Energy-based Feynman-Kac Corrector (EBFKC).}
Given a reference energy $E_t(\mathbf{x})$ defining the desired path $\pi_t \propto e^{-E_t(\mathbf{x})}$,
the paper derives a Feynman-Kac potential $G_t$ that, when incorporated as a
multiplicative weight in the SMC procedure, ensures the particle ensemble tracks
the prescribed path $\{p_t\}$ exactly in the continuous-time population limit,
regardless of the score function error $\epsilon_t$.

The key insight is that the Feynman-Kac formula from stochastic analysis provides a
closed-form expression for the correction term that neutralizes the effect of
$\epsilon_t$ without requiring access to the true score function---only the reference
energy is needed.

\section{Technical Formulation}

Let $\hat{s}_\theta(\mathbf{x},t)$ be the imperfect score with error
$\epsilon_t(\mathbf{x}) = \hat{s}_\theta - \nabla\log p_t$.
Standard SMC propagates under:
\[
d\mathbf{x}_t = \hat{s}_\theta(\mathbf{x}_t,t)\,dt + \sqrt{2}\,dW_t
\]
and tracks $p_t$ only when $\epsilon_t = 0$.

EBFKC derives the Feynman-Kac potential:
\[
G_t(\mathbf{x}) = \exp\!\left(-\int_0^t \bigl\langle \epsilon_\tau(\mathbf{x}_\tau),\,
  \nabla\log p_\tau(\mathbf{x}_\tau)\bigr\rangle \, d\tau\right)
\]
The corrected particle system with weights proportional to $G_t$ converges to $p_t$
in the population limit. A variance-controlling guidance term:
\[
\hat{s}_\theta^{\text{VC}}(\mathbf{x},t) = \hat{s}_\theta(\mathbf{x},t)
  - \gamma \nabla_{\mathbf{x}} \log w_t(\mathbf{x})
\]
prevents particle collapse by penalizing low-weight regions, where $\gamma > 0$
is a guidance strength parameter.

\section{Learning or Inference Procedure}

EBFKC requires no additional training---it operates on a frozen pretrained diffusion model.

\begin{enumerate}[itemsep=2pt,parsep=0pt]
  \item \textbf{Specify} reference energy $E_t$ (free energy functional or reward model).
  \item \textbf{Initialize} $N$ particles $\{\mathbf{x}_T^{(i)}\}_{i=1}^N$ from the
    prior $p_T$.
  \item \textbf{Propagate} with $\hat{s}_\theta$ plus variance-controlling guidance,
    accumulating Feynman-Kac weights $G_t^{(i)}$ at each step.
  \item \textbf{Resample} when the effective sample size
    $N_{\text{eff}} = (\sum_i w_i^2)^{-1}$ drops below $N/2$.
  \item \textbf{Return} the weighted ensemble $\{(\mathbf{x}_0^{(i)}, w_0^{(i)})\}$.
\end{enumerate}

The overhead relative to standard SMC is one energy evaluation per particle per step.

\section{What the Guarantee Says}

In the continuous-time population limit ($N \to \infty$, $\Delta t \to 0$), EBFKC
tracks the prescribed path $\{p_t\}$ exactly, independent of $\epsilon_t$.
For finite $N$, the error scales as $O(1/\sqrt{N})$ in particle count and
$O(\Delta t)$ in discretization step size. Crucially, unlike standard SMC,
the $N \to \infty$ limit converges to the correct target rather than a biased surrogate.

\section{Experimental Findings}

\begin{itemize}[itemsep=2pt,parsep=0pt]
  \item \textbf{Gaussian mixture models:} EBFKC recovers all modes of a multi-modal
    2D target; standard SMC with imperfect model overweights dominant modes.
  \item \textbf{Particle systems:} EBFKC accurately recovers equilibrium distributions
    for interacting particles; standard baselines show persistent density errors.
  \item \textbf{Alanine dipeptide:} EBFKC matches the ground-truth Ramachandran
    free-energy surface within statistical uncertainty; standard SMC exhibits errors
    of several $k_BT$ in high-barrier regions.
  \item \textbf{Alanine tetrapeptide:} EBFKC scales to the larger system and maintains
    accurate free-energy profiles under both annealing and reward-tilting regimes.
\end{itemize}

\section{Ablations and Interpretation}

\textbf{Particle count:} EBFKC converges faster to the correct target than standard
SMC as $N$ increases; at low $N$, variance-controlling guidance is essential.

\textbf{Model imperfection level:} As score error $\epsilon_t$ is artificially
increased, standard SMC degrades rapidly while EBFKC remains accurate.

\textbf{Variance-controlling guidance:} Removing this term causes particle collapse
at long horizons, confirming its necessity.

\textbf{Feynman-Kac vs.\ naive reweighting:} Simple importance weighting without
the Feynman-Kac path correction tracks a biased distribution even with large $N$,
underscoring that the correction must operate on the dynamics, not just the terminal
weights.

\section*{Reference}

Zuokai Wen, Louis Grenioux, Weinan E, Jiequn Han.
\textbf{Error-Corrected Inference-Time Scaling for Imperfect Diffusion Models}.
arXiv:2610.01933, October 2026.\\
\url{https://arxiv.org/abs/2610.01933}

\end{document}
